\chapter{Introduction}
\label{ch:intro}

\section{What is \qv?}

\qv{} is a research-grade optimiser for \emph{Ising} and \emph{QUBO} problems
--- the two canonical formulations into which an enormous range of hard
combinatorial optimisation problems (partitioning, scheduling, routing,
satisfiability, portfolio selection, protein folding, \dots) can be encoded.
It provides four annealing engines, spanning the spectrum from purely
classical thermal optimisation to the closest CPU-side simulation of a
physical quantum annealer:

\begin{center}
\small
\begin{tabular}{@{}llll@{}}
\toprule
\textbf{Method} & \textbf{Name} & \textbf{What it simulates} & \textbf{Ch.}\\
\midrule
Simulated Annealing & \code{sa} & Classical thermal fluctuations & \ref{ch:sa}\\
Simulated Quantum Annealing & \code{sqa} & Discrete-time path integral (Trotterised QA) & \ref{ch:sqa}\\
SQA + Parallel Tempering & \code{sqapt} & QA with replica exchange, $(\beta,\Gamma)$ ladder & \ref{ch:sqapt}\\
Continuous-Time PIMC & \code{ctpimc} & Continuous-time path integral (worldlines) & \ref{ch:ctpimc}\\
\bottomrule
\end{tabular}
\end{center}

All four engines share a single high-level Python entry point,
\code{solve()}, and a common C++17 core exposed through pybind11 bindings.
On top of the engines, \qv{} implements a theoretically grounded
\emph{optimal adaptive $\Jp$ schedule} (Chapter~\ref{ch:optimal}) derived
from the local adiabaticity condition of Roland and Cerf: the annealing
trajectory automatically slows down where the simulated system is close to
its quantum critical point and accelerates elsewhere.

\section{Design philosophy}

\begin{description}[style=nextline]
\item[One call to solve.] \code{solve(problem, method=..., ...)} accepts
NumPy arrays, dictionaries, edge lists, \code{dimod} binary quadratic models
and \code{networkx} graphs, auto-detects the input type, auto-tunes a
schedule from the problem's own energy scale, runs any number of independent
reads and returns the best result.

\item[Physics first.] Every knob in the library corresponds to a physical
quantity: $\beta$ is an inverse temperature, $\Gamma$ a transverse field,
$M$ a number of imaginary-time slices, $\Jp$ a Trotter bond strength.
The documentation (this manual) derives all of them so that parameter
choices can be reasoned about rather than guessed.

\item[Nothing hard-coded.] Schedule endpoints, temperature ladders and the
adaptive-schedule calibration all derive from the \emph{problem's} coupling
scale (via random-flip probing or the RMS coupling $J_{\mathrm{rms}}$),
so the same code works unchanged from $n=10$ toy instances to
$n>100\,000$ sparse graphs.

\item[Research instrumentation.] Observer classes capture per-sweep energy,
magnetisation, and full spin-state traces; the adaptive schedule exposes its
realised $\Jp$ trajectory and bond-susceptibility trace so that every run can
be audited.

\item[HPC ready.] OpenMP parallelism within a run, a multiprocessing/MPI
launcher across reads and nodes, and SLURM job-array recipes.
\end{description}

\section{Architecture at a glance}

\begin{figure}[ht]
\centering
\begin{tikzpicture}[
  layer/.style={rectangle,rounded corners=2pt,draw=qnavy!50,fill=qnavy!6,
                text width=13.2cm,align=center,minimum height=0.95cm,font=\small},
  lab/.style={font=\scriptsize\sffamily\color{qgray}},
  node distance=3.5mm]
\node[layer,fill=qteal!10,draw=qteal!60] (a)
  {\textbf{Python front end}\quad \code{solve()} \,\textbullet\, auto-tuned schedules \,\textbullet\, problem adapters (ndarray / dict / BQM / graph)};
\node[layer,below=of a] (b)
  {\textbf{pybind11 bindings}\quad models \,\textbullet\, schedules \,\textbullet\, annealers \,\textbullet\, observers \,\textbullet\, results};
\node[layer,below=of b] (c)
  {\textbf{C++17 core}\quad \code{Annealer} \,\textbullet\, \code{ReplicaAnnealer} \,\textbullet\, \code{SQAAnnealer} \,\textbullet\, \code{SQAParallelTemperingAnnealer} \,\textbullet\, \code{CTPIMCAnnealer}};
\node[layer,below=of c,fill=qnavy!12] (d)
  {\textbf{Parallel runtime}\quad OpenMP (replicas $\times$ slices) \,\textbullet\, optional MPI \,\textbullet\, SLURM array launcher};
\draw[-{Stealth},thick,qnavy!60] (a) -- (b);
\draw[-{Stealth},thick,qnavy!60] (b) -- (c);
\draw[-{Stealth},thick,qnavy!60] (c) -- (d);
\end{tikzpicture}
\caption{The \qv{} software stack. Users normally interact only with the top
layer; every lower layer remains individually scriptable.}
\label{fig:stack}
\end{figure}

\section{A sixty-second tour}

Install from the repository root and solve a three-variable QUBO:

\begin{bashcode}
python -m pip install . --no-build-isolation
\end{bashcode}

\begin{pycode}
import numpy as np
from qanneal import solve

# Minimise  E(x) = x0 + x2 - 2 x0 x1 - 2 x1 x2   over bits x in {0,1}^3
Q = np.array([
    [ 1.0, -2.0,  0.0],
    [-2.0,  0.0, -2.0],
    [ 0.0, -2.0,  1.0],
], dtype=float)

result = solve(Q, method="sqapt", reads=16, seed=0, return_bits=True)
print(result.best_sample)   # bit vector, e.g. [1 1 1]
print(result.best_energy)   # minimum QUBO energy found
\end{pycode}

Everything that just happened --- how the QUBO became an Ising Hamiltonian,
how a $(\beta,\Gamma)$ ladder was tuned to the matrix's energy scale, what a
Trotter slice is, why replica exchange helps, and how to make the run
adaptive, parallel or verifiable --- is explained in the rest of this manual.

\section{How to read this manual}

\begin{description}[style=nextline]
\item[Part I --- Foundations.] The Ising and QUBO models, their exact
equivalence, encodings of classic NP-hard problems, and the statistical
mechanics (Boltzmann distributions, Markov chains, detailed balance) on which
every engine rests. Read this if you want the mathematics from zero.

\item[Part II --- The annealing engines.] One chapter per engine. Each
derives the algorithm, states the exact acceptance rules implemented in the
C++ core, and ends with parameter guidance.

\item[Part III --- The optimal adaptive schedule.] The local-adiabaticity
derivation, the bond susceptibility $\chiB$, budget calibration, the
inverse-CDF trajectory, and how to diagnose a run.

\item[Part IV --- Using the library.] Installation, the full Python API
reference, worked end-to-end examples, parallel/HPC deployment,
troubleshooting, and the C++ API.
\end{description}

\begin{notebox}[Typographical conventions]
Python and C++ identifiers are set in \code{monospace}. Key equations are
boxed with a navy rule; physical intuition appears in plum ``Physical
picture'' boxes; pitfalls appear in amber warning boxes. Spins are
$s_i\in\{-1,+1\}$, bits are $x_i\in\{0,1\}$, $\beta$ is always an inverse
temperature and $M$ the number of Trotter slices.
\end{notebox}
